% ===== Джерела до розділу 16. Ключі з префіксом ch16: =====
% Позиції ядра (bib-core.tex) тут НЕ дублюються.

% --- Реагування на інциденти ---
\bibitem{ch16:nist-ir} Cichonski P., Millar T., Grance T., Scarfone K. Computer Security Incident Handling Guide. NIST Special Publication 800-61, Revision~2. Gaithersburg: NIST, 2012. DOI: 10.6028/NIST.SP.800-61r2.

\bibitem{ch16:sans-ir} Murray J. Incident Handler's Handbook. SANS Institute Reading Room, 2015.

\bibitem{ch16:isoir} International Organization for Standardization. ISO/IEC~27035: Information Technology --- Security Techniques --- Information Security Incident Management. Geneva: ISO, 2016--2023.

% --- Цифрова форензика ---
\bibitem{ch16:iso27037} International Organization for Standardization. ISO/IEC~27037: Guidelines for Identification, Collection, Acquisition and Preservation of Digital Evidence. Geneva: ISO, 2012.

\bibitem{ch16:casey} Casey E. Digital Evidence and Computer Crime: Forensic Science, Computers, and the Internet. 3rd ed. Waltham: Academic Press, 2011. 840~p.

\bibitem{ch16:carrier} Carrier B. File System Forensic Analysis. Boston: Addison-Wesley, 2005. 600~p.

\bibitem{ch16:orderofvolatility} Brezinski D., Killalea T. Guidelines for Evidence Collection and Archiving. RFC~3227. IETF, 2002.

% --- Відновлення та уроки ---
\bibitem{ch16:recovery} Souppaya M., Scarfone K. Guide for Cybersecurity Event Recovery. NIST Special Publication 800-184. Gaithersburg: NIST, 2016.

% --- Антифорензика ---
\bibitem{ch16:antiforensics} Garfinkel S. Anti-Forensics: Techniques, Detection and Countermeasures. \emph{Proceedings of the 2nd International Conference on i-Warfare and Security}. 2007. P.~77--84.

% --- Українські джерела ---
\bibitem{ch16:ir-ua} Організація процесу реагування на кіберінциденти в установах України відповідно до вимог законодавства. \emph{Кібербезпека: освіта, наука, техніка}. 2024. URL: \url{https://csecurity.kubg.edu.ua/index.php/journal/article/view/816}.

\bibitem{ch16:forensic-ua} Методи цифрової криміналістики при розслідуванні інцидентів, пов'язаних зі шкідливим програмним забезпеченням. \emph{Захист інформації}. 2023. Т.~25, №~2.
